% GENERATED FILE. Edit canonical lessons, then run npm run generate:books.
% canonical-source-hash: fnv1a64:d83f31ec744fad54
% canonical-lessons: ZH-C163-youkong, ZH-C163-pei, ZH-C163-yecan, ZH-C163-xiaci, ZH-C163-yikuair, ZH-R163-review-158-163

\chapter{Free This Weekend}
\label{ch:zh-free-this-weekend}
\hlchaptermodality{163}

\begin{chapteropening}
\textbf{By the end of this chapter:} \emph{I can ask if someone is free, go with them, plan a picnic and put it off until next time.}

Complete the last lesson of chapter 163: I can ask if someone is free, go with them, plan a picnic and put it off until next time.
\end{chapteropening}

\section[yǒukòng]{yǒukòng — to have free time; to be free}

\label{lesson:ZH-C163-youkong}



\begin{quote}
Before the new one: say the Mandarin for now and then, then the Mandarin for normally.
\end{quote}

\subsection*{You'll want to know: yǒukòng}
\emph{yǒukòng} — \textquotedblleft{}to have free time; to be free\textquotedblright{}.

\begin{quote}
\emph{nǐ zhōumò yǒukòng ma} — Are you free at the weekend?
\end{quote}

Its characters come later; for now it is a word for the ear.

\begin{cousinweb}
\emph{yǒukòng} is built from \zh{有} 'to have' and \emph{kòng} 'a gap', both inherited from Middle Chinese, \emph{hjuwX} and \emph{khuwngH} (Guangyun, in Baxter's transcription), which glosses the second 'a gap'. It is the falling-tone partner of \emph{kōng} 'empty': to have a gap in the day.
\end{cousinweb}

\begin{grammarlens}[title={What you've built}]
One more word for talking about plans.
\end{grammarlens}

\subsection*{Your turn}
\begin{itemize}
\raggedright
  \item \emph{Say it:} \emph{yǒukòng}
  \item \emph{Say it:} \emph{yǒukòng}, once more
  \item \emph{Say it:} \emph{nǐ zhōumò yǒukòng ma} — Are you free at the weekend?
  \item \emph{Recall:} say the Mandarin for now and then, then the Mandarin for normally, then say \emph{yǒukòng} again
\end{itemize}

\begin{tcolorbox}[breakable,colback=teal!4,colframe=teal!35!black,title={Before you move on}]
What does yǒukòng mean? (\textquotedblleft{}To have free time; to be free\textquotedblright{}.) Say it once more.
\end{tcolorbox}
\section[péi]{péi — to keep someone company; to go along with}

\label{lesson:ZH-C163-pei}



\begin{quote}
Before the new one: say the Mandarin for normally, then the Mandarin for to have free time.
\end{quote}

\subsection*{You'll want to know: péi}
\emph{péi} — \textquotedblleft{}to keep someone company; to go along with\textquotedblright{}.

\begin{quote}
\emph{zhōumò wǒ péi māma qù mǎi dōngxi} — At the weekend I'm going shopping with my mum.
\end{quote}

Its characters come later; for now it is a word for the ear.

\begin{cousinweb}
\emph{péi} is inherited from Middle Chinese \emph{bwoj} (Guangyun, in Baxter's transcription), glossed 'to attend at someone's side'.
\end{cousinweb}

\begin{grammarlens}[title={What you've built}]
One more word for talking about plans.
\end{grammarlens}

\subsection*{Your turn}
\begin{itemize}
\raggedright
  \item \emph{Say it:} \emph{péi}
  \item \emph{Say it:} \emph{péi}, once more
  \item \emph{Say it:} \emph{zhōumò wǒ péi māma qù mǎi dōngxi} — At the weekend I'm going shopping with my mum.
  \item \emph{Recall:} say the Mandarin for normally, then the Mandarin for to have free time, then say \emph{péi} again
\end{itemize}

\begin{tcolorbox}[breakable,colback=teal!4,colframe=teal!35!black,title={Before you move on}]
What does péi mean? (\textquotedblleft{}To keep someone company; to go along with\textquotedblright{}.) Say it once more.
\end{tcolorbox}
\section[yěcān]{yěcān — to have a picnic; a picnic}

\label{lesson:ZH-C163-yecan}



\begin{quote}
Before the new one: say the Mandarin for to have free time, then the Mandarin for to keep someone company.
\end{quote}

\subsection*{You'll want to know: yěcān}
\emph{yěcān} — \textquotedblleft{}to have a picnic; a picnic\textquotedblright{}.

\begin{quote}
\emph{zhōumò wǒmen qù gōngyuán yěcān} — This weekend we're going to the park for a picnic.
\end{quote}

Its characters come later; for now it is a word for the ear.

\begin{cousinweb}
\emph{yěcān} is built from \emph{yě} 'open country' and \emph{cān} 'meal', both inherited from Middle Chinese, \emph{yaeX} and \emph{tshan} (Guangyun, in Baxter's transcription); the dictionary glosses the second 'to swallow'. A meal out in the fields.
\end{cousinweb}

\begin{grammarlens}[title={What you've built}]
One more word for talking about plans.
\end{grammarlens}

\subsection*{Your turn}
\begin{itemize}
\raggedright
  \item \emph{Say it:} \emph{yěcān}
  \item \emph{Say it:} \emph{yěcān}, once more
  \item \emph{Say it:} \emph{zhōumò wǒmen qù gōngyuán yěcān} — This weekend we're going to the park for a picnic.
  \item \emph{Recall:} say the Mandarin for to have free time, then the Mandarin for to keep someone company, then say \emph{yěcān} again
\end{itemize}

\begin{tcolorbox}[breakable,colback=teal!4,colframe=teal!35!black,title={Before you move on}]
What does yěcān mean? (\textquotedblleft{}To have a picnic; a picnic\textquotedblright{}.) Say it once more.
\end{tcolorbox}
\section[xià cì]{xià cì — next time}

\label{lesson:ZH-C163-xiaci}



\begin{quote}
Before the new one: say the Mandarin for to keep someone company, then the Mandarin for to have a picnic.
\end{quote}

\subsection*{You'll want to know: xià cì}
\emph{xià cì} — \textquotedblleft{}next time\textquotedblright{}.

\begin{quote}
\emph{jīntiān wǒ méiyǒu kòng xià cì wǒmen yìqǐ qù} — I'm busy today; let's go together next time.
\end{quote}

Its characters come later; for now it is a word for the ear.

\begin{cousinweb}
\emph{xià cì} is built from \zh{下} 'below, next' and \emph{cì} 'a turn', both inherited from Middle Chinese, \emph{haeX} and \emph{tshijH} (Guangyun, in Baxter's transcription), which glosses the second 'order, sequence': the next turn.
\end{cousinweb}

\begin{grammarlens}[title={What you've built}]
One more word for talking about plans.
\end{grammarlens}

\subsection*{Your turn}
\begin{itemize}
\raggedright
  \item \emph{Say it:} \emph{xià cì}
  \item \emph{Say it:} \emph{xià cì}, once more
  \item \emph{Say it:} \emph{jīntiān wǒ méiyǒu kòng xià cì wǒmen yìqǐ qù} — I'm busy today; let's go together next time.
  \item \emph{Recall:} say the Mandarin for to keep someone company, then the Mandarin for to have a picnic, then say \emph{xià cì} again
\end{itemize}

\begin{tcolorbox}[breakable,colback=teal!4,colframe=teal!35!black,title={Before you move on}]
What does xià cì mean? (\textquotedblleft{}Next time\textquotedblright{}.) Say it once more.
\end{tcolorbox}
\section[yíkuàir]{yíkuàir — all together; together (in everyday speech)}

\label{lesson:ZH-C163-yikuair}



\begin{quote}
Before the new one: say the Mandarin for to have a picnic, then the Mandarin for next time.
\end{quote}

\subsection*{You'll want to know: yíkuàir}
\emph{yíkuàir} — \textquotedblleft{}all together; together (in everyday speech)\textquotedblright{}.

\begin{quote}
\emph{wǒmen yíkuàir qù yěcān} — Let's all go for the picnic together.
\end{quote}

Its characters come later; for now it is a word for the ear.

The tone of \zh{一} changes: yī is said yí before the fourth tone of kuài.

\begin{cousinweb}
\emph{yíkuàir} is built from \zh{一} 'one', \emph{kuài} 'a lump' and the ending \zh{儿}, all three inherited from Middle Chinese, \emph{'jit}, \emph{khwojH} and \emph{nye} (Guangyun, in Baxter's transcription). The dictionary glosses the second 'a clod of earth': in one lump.
\end{cousinweb}

\begin{grammarlens}[title={What you've built}]
One more word for talking about plans.
\end{grammarlens}

\subsection*{Your turn}
\begin{itemize}
\raggedright
  \item \emph{Say it:} \emph{yíkuàir}
  \item \emph{Say it:} \emph{yíkuàir}, once more
  \item \emph{Say it:} \emph{wǒmen yíkuàir qù yěcān} — Let's all go for the picnic together.
  \item \emph{Recall:} say the Mandarin for to have a picnic, then the Mandarin for next time, then say \emph{yíkuàir} again
\end{itemize}

\begin{tcolorbox}[breakable,colback=teal!4,colframe=teal!35!black,title={Before you move on}]
What does yíkuàir mean? (\textquotedblleft{}All together; together (in everyday speech)\textquotedblright{}.) Say it once more.
\end{tcolorbox}
\section[(dialogue)]{Looking back: from Waking Up to Free This Weekend}

\label{lesson:ZH-R163-review-158-163}



\begin{quote}
Say the Mandarin for to wake someone up and to fall asleep.
\end{quote}

\subsection*{Your turn}
\begin{itemize}
\raggedright
  \item \emph{Say it:} to wake someone up, to fall asleep, to have a dream, to stand up, to be used to
  \item \emph{Say it:} a bedroom, a living room, a study at home, a balcony, a window
  \item \emph{Say it:} to sweep the floor, to do the washing-up, to make the bed, to get the jobs done, to look after
  \item \emph{Say it:} uncooked rice, measure word for a serving, fried rice, beef, chicken meat
  \item \emph{Say it:} as a rule, every single time, regularly, now and then, normally
  \item \emph{Say it:} to have free time, to keep someone company, to have a picnic, next time, all together
  \item \emph{Read it:} \zh{房} \zh{米} — each one aloud, then a word that holds it
\end{itemize}

\begin{tcolorbox}[breakable,colback=teal!4,colframe=teal!35!black,title={Before you move on}]
To wake someone up? (\emph{jiàoxǐng}.) To fall asleep? (\emph{shuìzháo}.) To have a dream? (\emph{zuòmèng}.) To stand up? (\textbf{\zh{起来}}, \emph{qǐlai}.) To be used to? (\emph{xíguàn}.) A bedroom? (\emph{wòshì}.) A living room? (\emph{kètīng}.) A study at home? (\textbf{\zh{书房}}, \emph{shūfáng}.) A balcony? (\emph{yángtái}.) A window? (\emph{chuānghu}.) To sweep the floor? (\emph{sǎodì}.) To do the washing-up? (\emph{xǐwǎn}.) To make the bed? (\emph{pūchuáng}.) To get the jobs done? (\emph{gànhuó}.) To look after? (\emph{zhàogù}.) Uncooked rice? (\textbf{\zh{大米}}, \emph{dàmǐ}.) Measure word for a serving? (\emph{fèn}.) Fried rice? (\emph{chǎofàn}.) Beef? (\emph{niúròu}.) Chicken meat? (\emph{jīròu}.) As a rule? (\emph{yìbān}.) Every single time? (\emph{zǒngshì}.) Regularly? (\emph{jīngcháng}.) Now and then? (\emph{ǒu'ěr}.) Normally? (\emph{tōngcháng}.) To have free time? (\emph{yǒukòng}.) To keep someone company? (\emph{péi}.) To have a picnic? (\emph{yěcān}.) Next time? (\emph{xià cì}.) All together? (\emph{yíkuàir}.)
\end{tcolorbox}
